
This dissertation presents a study of neutral meson production in proton-proton ($p+p$) collisions at $\sqrt{s}$ = 200 \ac{GeV} collected by newly built \ac{sPHENIX} detector. \ac{sPHENIX} detector, within \ac{RHIC} facility, is a state-of-the-art successor to \ac{PHENIX} with extended pseudorapidity ($\eta$) and complete $\phi$-coverage. Our research focuses on the measurement of invariant cross section of $\eta$ meson and $\pi^{0}$ via their two-photon decay mode. The final results are compared with results from \ac{PHENIX} experiment to provide the initial validation of data measured by \ac{sPHENIX} detector. In particular, a first full physics measurement of the ratio of $\eta$ meson to $\pi^0$ production cross sections is made and found to be in good agreement with past \ac{PHENIX} measurement. 

Furthermore, this dissertation reports various background reconstruction strategies and the application of supervised machine learning to isolate $\eta$ meson signals from background. Ultimately, we discuss the implementation of calibration methods of \ac{sPHENIX} \ac{EMCal} using $\pi^0$ decaying to two photons, as well as the methods' development and validation.
% <-- IMPORTANT: keep this as the last line and no space above this line to keep everything in frontmatter otherwise we might get an error